\justifying
\NSFCBodyText

\subsubsection{项目特色}
本项目把矿山变化检测中的误报和细小变化漏检放在同一对照体系中研究，以一致性约束的作用区域为切入点，将方法效果追溯到约束对象、变化尺度和外观差异。研究使用公开数据与矿区隔离协议，保留失效矿区和反例，强调条件可重复，而非只追求单个总体指标。

可靠性评价与检测方法采用分步冻结：先选择检测器，再检验其错误排序，把概率误差、复核面积和全图剩余错误分别记录。这样可以判断收益究竟来自检测改进还是复核选择，并明确高置信度漏检仍然存在的范围。

\subsubsection{创新点}
\textbf{（1）检验一致性约束空间分配的条件性作用。}
已有特征约束和域泛化方法构成本项目的比较基础。拟通过固定面积、随机掩码分配与近邻方法局部替换，将背景约束的性能收益与约束误覆盖变化区域时的检出损失分开，识别其受变化尺度和外观差异影响的适用边界。创新成立的证据应同时包含近邻增量与可排除替代解释的对照；若收益由平滑或采样难度解释，则收缩为失败条件研究，不把掩码替换本身称为新机制。

\textbf{（2）围绕全图剩余错误建立复核效用证据。}
将校准与推理分歧置于同一冻结检测器、同一矿区面积预算下比较，统一覆盖预测变化和未变化区域，分别检验误报与漏检的残留。预期贡献是明确可靠性排序在矿山变化复核中的有效条件及盲区，所用温度缩放、熵与选择性预测均为已有工具；离线结果仅表述为理想纠错上限。
